%% kapa: exam=1 round=25 year=2014 period=1 form=A subject=civil_law subject_ko=민법 seq=22 number=22
%% source: Q-Net 감정평가사 기출문제 (cst00309) artlSeq=5101207 file=extracted/1008/5101207/2014년 제25회 감정평가사 1차시험 기출문제/1교시(A형).hwp
%% conversion: kordoc 4.24.51 -> tools/gen1.py (auto); review=proofread-sonnet; flags=none
\begin{question}{22}
표현대리에 관한 설명으로 옳지 않은 것은? (다툼이 있으면 판례에 의함)
\begin{choices}
  \item 권한을 넘은 표현대리인지를 판단할 때 정당한 이유의 유무는 대리행위 시를 기준으로 한다.
  \item 권한을 넘은 표현대리에서는 기본대리권의 내용이 되는 행위와 표현대리 행위가 반드시 같은 종류의 것이어야 한다.
  \item 부부간의 일상가사대리권을 기본대리권으로 하여 권한을 넘은 표현대리가 성립할 수 있다.
  \item 복대리인의 대리권은 권한을 넘은 표현대리의 기본대리권이 될 수 있다.
  \item 표현대리가 성립하여도 무권대리의 성질이 유권대리로 전환되는 것은 아니다.
\end{choices}
\end{question}
\answer{2}
